
\begin{table}[t]
\centering
\caption{Dataset Description and Experimental Splits}
\label{tab:dataset}
\small
\begin{tabular}{@{}llrl@{}}
\toprule
\textbf{Split} & \textbf{Source} & \textbf{Size} & \textbf{Purpose} \\
\midrule
Warmup          & RouteLLM Battles & 80,000 & PCA training (384$\rightarrow$32) + LinUCB priors ($\mathbf{A}$, $\mathbf{b}$) \\
Development     & LMSYS Arena      & 1,121 & Online learning \& calibration \\
Holdout         & LMSYS Arena      & 750 & Final evaluation \\
\midrule
\textbf{Total} & & \textbf{81,871} & \\
\bottomrule
\end{tabular}

\vspace{1em}
\footnotesize
\textbf{Data Sources and Independence.} All prompts originate from the LMSYS Chat Arena ecosystem~\cite{zheng2023lmsys}, but the warmup and evaluation datasets are \emph{independent collections} from different data sources, sampling periods, and prompt populations.
\textbf{RouteLLM Battles}~\cite{ong2024routellm}: 80K pairwise comparisons (mixtral-8x7b-instruct vs gpt-4-turbo) from the HuggingFace dataset \texttt{routellm/gpt4\_judge\_battles}---a curated battle collection. Used for PCA training (384$\rightarrow$32) and warmup prior generation ($\mathbf{A} \in \mathbb{R}^{33 \times 33}$, $\mathbf{b} \in \mathbb{R}^{33}$).
\textbf{LMSYS Arena (evaluation):} Dev and holdout prompts from the LMSYS general prompt pool, collected independently of the RouteLLM battles. Same model pair but different source, period, and prompt population. This independence ensures that the PCA has never seen the evaluation prompts---no decontamination step is needed; the datasets are disjoint by provenance.
\textbf{Data Quality.} Zero data leakage verified via automated checks (243 incidentally overlapping prompts removed, 0.24\%; overlap is due to both datasets sampling from the broader LMSYS user base, not shared provenance). Dev and holdout sets created using stratified sampling across three dimensions: category (STEM, CODE, GENERAL), complexity (Low, Med, High), and difficulty (Easy, Hard, Contentious), resulting in representative coverage across diverse prompt characteristics.
\textbf{Reward Structure.} Rewards are discrete pairwise preference outcomes (win/tie/loss) from LMSYS Chatbot Arena human evaluations, consistent with the categorical analysis in Figure~\ref{fig:lmsys_holdout_structure}. Of 750 holdout prompts, 72.8\% are ties (routing-irrelevant); only 204 prompts (27.2\%) have differential model performance.
\textbf{Sample Size.} Development set (1,121 prompts) enables online learning with sufficient data for bandit convergence. Holdout set (750 prompts) determined by data availability. Monte-Carlo power analysis (Appendix~\ref{app:power}), consistent with Figure~\ref{fig:lmsys_holdout_structure}'s categorical approach, uses McNemar's exact test (paired binary data) and binomial test (routing accuracy on informative prompts): 80\% power to detect routing accuracy $\geq$58\% on the 204 informative prompts ($\alpha = 0.05$). Multi-seed validation (Table~\ref{tab:corralling}, $N$=10--30 seeds) provides additional robustness evidence.
\end{table}
